% GENERATED FILE. Edit canonical lessons, then run npm run generate:books.
% canonical-source-hash: fnv1a64:273b74efc2b4e62d
% canonical-lessons: SA-C419-vapi, SA-C419-utsah, SA-C419-puskarini, SA-C419-kulya, SA-C419-palvalam

\chapter{Reservoirs, Springs and Pools}
\label{ch:sa-reservoirs-springs-and-pools}
\hlchaptermodality{419}

\begin{chapteropening}
\textbf{By the end of this chapter:} \emph{I can say where to find water on a journey.}

Complete the last lesson of chapter 419: I can say where to find water on a journey.
\end{chapteropening}

\section[\texorpdfstring{\sk{वापी} (vāpī)}{वापी (vāpī)}]{\sk{वापी} (vāpī) — an oblong reservoir; a tank}

\label{lesson:SA-C419-vapi}



\begin{quote}
Before the new one: say the Sanskrit for a pearl, then the Sanskrit for coral.
\end{quote}

\subsection*{You'll want to know: \sk{वापी}}
\textbf{\sk{वापी}} — \emph{vāpī} — \textquotedblleft{}an oblong reservoir; a tank\textquotedblright{}.

\begin{quote}
\sk{ग्रामस्य} \sk{मध्ये} \sk{गभीरा} \sk{वापी} \sk{अस्ति।}
\end{quote}

\begin{quote}
\emph{grāmasya madhye gabhīrā vāpī asti.}
\end{quote}

\begin{quote}
In the middle of the village there is a deep reservoir.
\end{quote}

A feminine noun.

\begin{cousinweb}
\textbf{\sk{वापी}} (\emph{vāpī}) comes from the root √vap, \textquotedblleft{}to strew, to heap up\textquotedblright{}, which also gives \emph{vapati}, \textquotedblleft{}sows\textquotedblright{}: MW gives \emph{vāpī} as a pond made by damming up earth, an oblong reservoir, a tank; it lives on as Hindi \emph{bāuṛī}, \textquotedblleft{}a well with steps\textquotedblright{} (MW; CDIAL 11529).
\end{cousinweb}

\begin{grammarlens}[title={What you've built}]
One more word for getting around.
\end{grammarlens}

\subsection*{Your turn}
\begin{itemize}
\raggedright
  \item \emph{Say it:} \emph{vāpī}
  \item \emph{Say it:} \emph{vāpī}, once more
  \item \emph{Say it:} \emph{grāmasya madhye gabhīrā vāpī asti.} — In the middle of the village there is a deep reservoir.
  \item \emph{Recall:} say the Sanskrit for a pearl, then the Sanskrit for coral, then say \emph{vāpī} again
\end{itemize}

\begin{tcolorbox}[breakable,colback=teal!4,colframe=teal!35!black,title={Before you move on}]
What does \sk{वापी} mean? (\textquotedblleft{}An oblong reservoir; a tank\textquotedblright{}.) Say it once more.
\end{tcolorbox}
\section[\texorpdfstring{\sk{उत्सः} (utsaḥ)}{उत्सः (utsaḥ)}]{\sk{उत्सः} (utsaḥ) — a spring of water}

\label{lesson:SA-C419-utsah}



\begin{quote}
Before the new one: say the Sanskrit for coral, then the Sanskrit for an oblong reservoir.
\end{quote}

\subsection*{You'll want to know: \sk{उत्सः}}
\textbf{\sk{उत्सः}} — \emph{utsaḥ} — \textquotedblleft{}a spring of water\textquotedblright{}.

\begin{quote}
\sk{पर्वते} \sk{शीतलः} \sk{उत्सः} \sk{अस्ति।}
\end{quote}

\begin{quote}
\emph{parvate śītalaḥ utsaḥ asti.}
\end{quote}

\begin{quote}
There is a cool spring on the mountain.
\end{quote}

A masculine noun.

\begin{cousinweb}
\textbf{\sk{उत्सः}} (\emph{utsaḥ}) comes from the root √ud, \textquotedblleft{}to wet\textquotedblright{}, as MW gives it; MW gives \emph{utsa} as a spring or fountain (MW).
\end{cousinweb}

\begin{grammarlens}[title={What you've built}]
One more word for getting around.
\end{grammarlens}

\subsection*{Your turn}
\begin{itemize}
\raggedright
  \item \emph{Say it:} \emph{utsaḥ}
  \item \emph{Say it:} \emph{utsaḥ}, once more
  \item \emph{Say it:} \emph{parvate śītalaḥ utsaḥ asti.} — There is a cool spring on the mountain.
  \item \emph{Recall:} say the Sanskrit for coral, then the Sanskrit for an oblong reservoir, then say \emph{utsaḥ} again
\end{itemize}

\begin{tcolorbox}[breakable,colback=teal!4,colframe=teal!35!black,title={Before you move on}]
What does \sk{उत्सः} mean? (\textquotedblleft{}A spring of water\textquotedblright{}.) Say it once more.
\end{tcolorbox}
\section[\texorpdfstring{\sk{पुष्करिणी} (puṣkariṇī)}{पुष्करिणी (puṣkariṇī)}]{\sk{पुष्करिणी} (puṣkariṇī) — a lotus pool}

\label{lesson:SA-C419-puskarini}



\begin{quote}
Before the new one: say the Sanskrit for an oblong reservoir, then the Sanskrit for a spring of water.
\end{quote}

\subsection*{You'll want to know: \sk{पुष्करिणी}}
\textbf{\sk{पुष्करिणी}} — \emph{puṣkariṇī} — \textquotedblleft{}a lotus pool\textquotedblright{}.

\begin{quote}
\sk{पुष्करिण्यां} \sk{पद्मानि} \sk{सन्ति।}
\end{quote}

\begin{quote}
\emph{puṣkariṇyāṃ padmāni santi.}
\end{quote}

\begin{quote}
There are lotuses in the lotus pool.
\end{quote}

A feminine noun.

\begin{cousinweb}
\textbf{\sk{पुष्करिणी}} (\emph{puṣkariṇī}) is built on \emph{puṣkara}, \textquotedblleft{}a blue lotus\textquotedblright{}, with the ending \emph{-iṇī}; MW gives \emph{puṣkariṇī} as a lotus pool, any pool or pond (MW).
\end{cousinweb}

\begin{grammarlens}[title={What you've built}]
One more word for getting around.
\end{grammarlens}

\subsection*{Your turn}
\begin{itemize}
\raggedright
  \item \emph{Say it:} \emph{puṣkariṇī}
  \item \emph{Say it:} \emph{puṣkariṇī}, once more
  \item \emph{Say it:} \emph{puṣkariṇyāṃ padmāni santi.} — There are lotuses in the lotus pool.
  \item \emph{Recall:} say the Sanskrit for an oblong reservoir, then the Sanskrit for a spring of water, then say \emph{puṣkariṇī} again
\end{itemize}

\begin{tcolorbox}[breakable,colback=teal!4,colframe=teal!35!black,title={Before you move on}]
What does \sk{पुष्करिणी} mean? (\textquotedblleft{}A lotus pool\textquotedblright{}.) Say it once more.
\end{tcolorbox}
\section[\texorpdfstring{\sk{कुल्या} (kulyā)}{कुल्या (kulyā)}]{\sk{कुल्या} (kulyā) — a water channel; a canal}

\label{lesson:SA-C419-kulya}



\begin{quote}
Before the new one: say the Sanskrit for a spring of water, then the Sanskrit for a lotus pool.
\end{quote}

\subsection*{You'll want to know: \sk{कुल्या}}
\textbf{\sk{कुल्या}} — \emph{kulyā} — \textquotedblleft{}a water channel; a canal\textquotedblright{}.

\begin{quote}
\sk{कुल्या} \sk{क्षेत्रेषु} \sk{जलं} \sk{वहति।}
\end{quote}

\begin{quote}
\emph{kulyā kṣetreṣu jalaṃ vahati.}
\end{quote}

\begin{quote}
The channel carries water to the fields.
\end{quote}

A feminine noun.

\begin{cousinweb}
\textbf{\sk{कुल्या}} (\emph{kulyā}) is the word MW gives for a small river, a canal or a channel for water; MW gives no derivation, so its origin is uncertain (MW).
\end{cousinweb}

\begin{grammarlens}[title={What you've built}]
One more word for getting around.
\end{grammarlens}

\subsection*{Your turn}
\begin{itemize}
\raggedright
  \item \emph{Say it:} \emph{kulyā}
  \item \emph{Say it:} \emph{kulyā}, once more
  \item \emph{Say it:} \emph{kulyā kṣetreṣu jalaṃ vahati.} — The channel carries water to the fields.
  \item \emph{Recall:} say the Sanskrit for a spring of water, then the Sanskrit for a lotus pool, then say \emph{kulyā} again
\end{itemize}

\begin{tcolorbox}[breakable,colback=teal!4,colframe=teal!35!black,title={Before you move on}]
What does \sk{कुल्या} mean? (\textquotedblleft{}A water channel; a canal\textquotedblright{}.) Say it once more.
\end{tcolorbox}
\section[\texorpdfstring{\sk{पल्वलम्} (palvalam)}{पल्वलम् (palvalam)}]{\sk{पल्वलम्} (palvalam) — a puddle; a small pond}

\label{lesson:SA-C419-palvalam}



\begin{quote}
Before the new one: say the Sanskrit for a lotus pool, then the Sanskrit for a water channel.
\end{quote}

\subsection*{You'll want to know: \sk{पल्वलम्}}
\textbf{\sk{पल्वलम्}} — \emph{palvalam} — \textquotedblleft{}a puddle; a small pond\textquotedblright{}.

\begin{quote}
\sk{वर्षायाः} \sk{अनन्तरं} \sk{मार्गे} \sk{पल्वलानि} \sk{सन्ति।}
\end{quote}

\begin{quote}
\emph{varṣāyāḥ anantaraṃ mārge palvalāni santi.}
\end{quote}

\begin{quote}
After the rain there are puddles on the road.
\end{quote}

A neuter noun.

\begin{cousinweb}
\textbf{\sk{पल्वलम्}} (\emph{palvalam}) is the word MW gives for a pool or small pond; MW compares Greek \emph{pēlos}, \textquotedblleft{}mud\textquotedblright{}, and Latin \emph{palus}, \textquotedblleft{}a marsh\textquotedblright{}, with a query, so its origin is uncertain (MW).
\end{cousinweb}

\begin{grammarlens}[title={What you've built}]
One more word for getting around.
\end{grammarlens}

\subsection*{Your turn}
\begin{itemize}
\raggedright
  \item \emph{Say it:} \emph{palvalam}
  \item \emph{Say it:} \emph{palvalam}, once more
  \item \emph{Say it:} \emph{varṣāyāḥ anantaraṃ mārge palvalāni santi.} — After the rain there are puddles on the road.
  \item \emph{Recall:} say the Sanskrit for a lotus pool, then the Sanskrit for a water channel, then say \emph{palvalam} again
\end{itemize}

\begin{tcolorbox}[breakable,colback=teal!4,colframe=teal!35!black,title={Before you move on}]
What does \sk{पल्वलम्} mean? (\textquotedblleft{}A puddle; a small pond\textquotedblright{}.) Say it once more.
\end{tcolorbox}
